\begin{helper}
Let \(m:\mathbb N\to\mathbb N\) satisfy
\[
  m(n)=\noderef{TwoBiteNaturalReducedVertexCount}(n)
\]
for every \(n\), and put
\[
  p(n)=\noderef{TwoBiteEdgeProbability}(1/2,n).
\]
There is a constant \(c>0\) such that, for all sufficiently large \(n\),
\[
\noderef{TwoBiteEventProbability}(n,m(n),p(n),\mathcal B_n)
  \le m(n)^2\exp(-c(\log n)^3),
\]
where \(\mathcal B_n\) is the event that both of the following hold for the
sample \(\omega\):

\[
\forall x\in\noderef{TwoBiteBaseVertex}(m(n)),\qquad
  |\noderef{TwoBiteLiftedNeighborhood}(\omega,x)|\le 2p(n)n,
\]
and there is an ordered blue pair \(b,c\in\operatorname{Fin}(m(n))\) with
\(b\ne c\) such that
\[
\begin{aligned}
&\left|
  \noderef{TwoBiteRedProjection}(\omega)
      [\noderef{TwoBiteLiftedNeighborhood}(\omega,\operatorname{Sum.inr} b)]
  \cap
  \noderef{TwoBiteRedProjection}(\omega)
      [\noderef{TwoBiteLiftedNeighborhood}(\omega,\operatorname{Sum.inr} c)]
 \right|  \\
&\qquad >
\left|
  \noderef{TwoBiteLiftedNeighborhood}(\omega,\operatorname{Sum.inr} b)
  \cap
  \noderef{TwoBiteLiftedNeighborhood}(\omega,\operatorname{Sum.inr} c)
\right|
  +100(\log n)^3 .
\end{aligned}
\]
\end{helper}
\begin{proof}
Fix \(m\) satisfying the displayed identity, and write
\[
  M=m(n),\qquad p=p(n),\qquad L=\log n .
\]
All estimates below are asserted only after increasing \(n\) past a fixed
threshold.  In particular \(L>0\), \(p\in[0,1]\) by
\(\noderef{OppositeProjectionEdgeProbBounds}\), and
\[
M=\noderef{TwoBiteNaturalReducedVertexCount}(n)
  =\left\lfloor\frac{n}{L^2}\right\rfloor .
\]
For each ordered blue pair \((b,c)\) with \(b\ne c\), let \(E_{b,c}\) be the
event that the global size bound holds and this pair violates the displayed
projected-intersection estimate.  The event \(\mathcal B_n\) is contained in
the union of the events \(E_{b,c}\), and there are at most \(M^2\) ordered
blue pairs.

We bound one fixed pair uniformly.  Apply
\(\noderef{OppositeProjectionBlueConditioning}\) to condition on the blue
base graph \(G_B\) and the blue coordinate map \(\rho:\operatorname{Fin} n
\to \operatorname{Fin} M\).  Cells with zero total mass contribute zero by
the conditional-probability convention.  On a nonzero cell, the lifted blue
neighborhoods of \(b\) and \(c\) are fixed sets
\[
  U=\noderef{TwoBiteLiftedNeighborhood}(\omega,\operatorname{Sum.inr} b),
  \qquad
  V=\noderef{TwoBiteLiftedNeighborhood}(\omega,\operatorname{Sum.inr} c).
\]
Let \(U_0=U\setminus V\) and \(V_0=V\setminus U\).  The size bound gives
\(|U_0|\le 2pn\) and \(|V_0|\le 2pn\).

The only projected intersections beyond the true intersection \(U\cap V\)
are vertices of \(V_0\) whose red coordinate lies in the red projection of
\(U_0\); this is exactly the containment supplied by
\(\noderef{OppositeProjectionIntersectionContainment}\).  After exposing the
red coordinates of \(U_0\), write the exposure as \(\eta\).  Then
\(|\pi_R(U_0)|\le |U_0|\), and the row-disjointness hypotheses required by
\(\noderef{OppositeProjectionBlueExposureTailBound}\) hold for \(U_0\) and
\(V_0\).  Hence, for \(t=\lceil100(\log n)^3\rceil\), each exposure cell has
conditional probability at most
\[
  \binom{|V_0|}{t}\left(\frac{|U_0|}{m(n)}\right)^t.
\]
Averaging over the exposure cells preserves the same bound, since the
right-hand side is independent of the exposure and zero-mass cells contribute
zero.

The containment shows that if \(E_{b,c}\) occurs, then the number \(W\) of
vertices of \(V_0\) whose red coordinate lies in the red projection of
\(U_0\) is strictly larger than \(100L^3\).  Since \(W\) is an integer, this
implies \(W\ge t\) for
\[
  t=\left\lceil100L^3\right\rceil .
\]
Using \(\binom bt\le (eb/t)^t\), the fixed-pair failure probability in a
blue conditioning cell is therefore at most
\[
  \Pr(E_{b,c}\mid G_B,\rho)
    \le \left(\frac{e|U_0||V_0|}{Mt}\right)^t .
\]
If \(t>|V_0|\), then the tail event is empty, so the same upper bound is
valid trivially.

It remains to estimate the parameter in this last expression.  Since
\[
p=\frac12\sqrt{\frac{L}{n}},
\]
we have \(p^2n=L/4\).  Also \(n/L^2\to\infty\), so eventually
\(n/L^2\ge2\).  For every real \(x\ge2\), \(\lfloor x\rfloor\ge x-1\ge x/2\);
therefore eventually
\[
  M=\left\lfloor\frac{n}{L^2}\right\rfloor
    \ge \frac{n}{2L^2}.
\]
Combining this lower bound for \(M\) with \(|U_0|,|V_0|\le2pn\) gives
\[
  \frac{|U_0||V_0|}{M}
    \le \frac{(2pn)^2}{M}
    =\frac{4p^2n^2}{M}
    =\frac{Ln}{M}
    \le 2L^3 .
\]
Because \(t=\lceil100L^3\rceil\ge100L^3\), we obtain
\[
  \frac{e|U_0||V_0|}{Mt}
    \le \frac{2eL^3}{100L^3}
    =\frac e{50}<1 .
\]
Set
\[
  c=100\log(50/e).
\]
This constant is positive, since \(50/e>1\).  The preceding inequality and
the monotonicity of \(a^s\) in \(s\) for \(0<a<1\) give
\[
\left(\frac{e|U_0||V_0|}{Mt}\right)^t
  \le \left(\frac e{50}\right)^t
  \le \left(\frac e{50}\right)^{100L^3}
  = \exp(-cL^3).
\]
Thus every fixed ordered blue pair has conditional failure probability at
most \(\exp(-cL^3)\), uniformly in the blue conditioning data.  Averaging over
the blue conditioning cells removes the conditioning and preserves this
uniform bound, with zero-mass cells contributing zero.  Finally, the union
bound over the at most \(M^2=m(n)^2\) ordered blue pairs gives
\[
  \Pr(\text{blue bad event})
    \le m(n)^2\exp(-c(\log n)^3),
\]
as claimed.
\end{proof}
